\documentclass[11pt]{article}
\usepackage[margin=1in]{geometry}
\usepackage{amsmath,amssymb,amsthm}
\usepackage[T1]{fontenc}
\usepackage{lmodern}
\usepackage[hidelinks]{hyperref}
\setlength{\emergencystretch}{1em}
\newtheorem{theorem}{Theorem}
\newtheorem{lemma}{Lemma}
\newcommand{\sd}{\operatorname{sd}}
\newcommand{\D}{\Delta}
\title{Flagness and Chain-Count Invariants\newline
  of Combinatorial Barycentric Subdivision}
\author{idealistichacker\\\small AI-assisted submission for conjecture 00000007126}
\date{October 23, 2026}
\begin{document}
\maketitle
\begin{abstract}
For every partially ordered set, its nonempty finite chains form an actual
abstract simplicial complex. We prove this complex is flag, its
$d$-dimensional face count is the number of chains with $d+1$ elements, and
these counts are invariant under order isomorphisms. Applying this construction
to a polytope's nonempty face poset is the standard \emph{combinatorial}
barycentric subdivision. We explicitly distinguish the whole polytope from its
boundary and use actual face subsets, without a simplicial-input restriction.
The source's unspecified invariant is interpreted as the enumerative
$f$-vector, not a complete classifier or a quantity preserved by performing
subdivision. No formal geometric triangulation theorem is claimed.
\end{abstract}

\section{Source Scope and Conventions}
The exact English source asserts the flag property of barycentric subdivision
and describes its invariant as chain counting of the original faces. The
Chinese source similarly refers to its combinatorial type, flagness and
face-chain counts. Neither defines a particular invariant or says that numeric
counts alone classify the subdivision. We prove the explicitly declared
standard combinatorial interpretation: the entire chain-incidence family
describes the abstract subdivision, and its $f$-vector is obtained by counting
chains by their number of elements. This vector is invariant under
combinatorial isomorphism; it is not asserted unchanged by subdivision.

The domain is nonempty bounded convex polytopes. For a polytope $Q$, let
$\mathcal F(Q)$ be its actual finite face family, with faces regarded as
\emph{subsets} of the ambient affine space. Define
\[
 P_{\rm full}(Q)=\{F\in\mathcal F(Q):F\ne\varnothing\},\qquad
 P_{\partial}(Q)=\{F\in\mathcal F(Q):F\ne\varnothing,\ F\ne Q\},
\]
both ordered by inclusion. The full convention includes $Q$; the boundary
convention does not. Empty faces are not subdivision vertices. In particular,
the boundary of a point has no nonempty faces. The combinatorial subdivision
is \emph{defined} to be the order complex of the corresponding face poset.
This standard description also applies to nonsimplicial polytopes
\cite[p.~11]{BW}; more generally, the face-poset order complex of a polyhedral
complex gives its combinatorial barycentric subdivision \cite[Section~2]{OS}.

Frozen source provenance:
\begin{quote}\small\ttfamily
Commit: b355e42906e6eaf52bc7e9d005e0ef02d8aa7400\\
Blob: 6c54d99d7caf701ff7e55c4530c2eb441a049aec.
\end{quote}
Both complete source languages are retained in the package's source document.

\section{The Actual Order Complex is Flag}
Let $P$ be any partially ordered set, not necessarily finite. A strict finite
chain means a finite subset totally ordered by $\le$. Repetitions are absent
because it is a set. Chains need not be saturated and may skip ranks. The
length of a chain means its number of elements, not its number of inequalities.
Define the abstract complex
\[
 \D(P)=\{s\subseteq P:s\text{ finite, nonempty, and a chain}\}.
\]
Every nonempty subset of a chain is a chain, and every singleton is a chain.
These facts prove the actual Mathlib abstract-complex structure fields.
The empty set is not a stored face under Mathlib's convention.

The comparability graph has vertex set $P$ and adjacency
\[
 x\sim y\quad\Longleftrightarrow\quad x\ne y\ \text{ and }\ (x\le y\text{ or }y\le x).
\]
It is an actual simple graph: adjacency is symmetric and irreflexive. Its
edges are exactly the two-element faces of $\D(P)$. A finite set is a clique
if and only if it is a chain, and therefore
\[
 s\in\D(P)\quad\Longleftrightarrow\quad s\ne\varnothing
       \ \text{ and }\ s\text{ is a clique in the comparability graph}.
\]

\begin{theorem}
For every poset $P$, the abstract order complex $\D(P)$ is flag.
\end{theorem}
\begin{proof}
Let $s$ be a nonempty finite vertex set such that every distinct pair in $s$
is an edge. Such pairs are comparable. Equal pairs are comparable by reflexivity.
Thus $s$ is a chain and is a face of $\D(P)$. This proves the clique-to-face
condition, with no assumed flagness or closure predicate.
\end{proof}

\section{Chain Counts and the Correct Invariant}
Assume $P$ is finite. Put
\[
 \mathcal S_k(P)=\{s\in\D(P):|s|=k\},\qquad
 \mathcal C_k(P)=\{s\subseteq P:s\ne\varnothing,\ s\text{ a chain},\ |s|=k\}.
\]
Membership in the actual complex gives an actual bijection
$\mathcal S_k(P)\simeq\mathcal C_k(P)$, by sending the finite vertex set to
the same underlying finite chain, with inverse inclusion as a face. Both
directions are proved; no counting identity is an external hypothesis.
By definition a $d$-dimensional abstract simplex has $d+1$ vertices, so
\[
 f_d(\D(P))=|\mathcal S_{d+1}(P)|=|\mathcal C_{d+1}(P)|.
\]
For empty $P$, all these nonempty-face counts are zero.

\begin{theorem}
If $e:P\simeq Q$ is an order isomorphism of finite posets, then
$f_d(\D(P))=f_d(\D(Q))$ for every $d\ge0$.
\end{theorem}
\begin{proof}
An order isomorphism preserves and reflects comparability, so the finite image
$e(s)$ is a chain if and only if $s$ is. Injectivity preserves its cardinality,
and nonemptiness is preserved. Taking finite images under $e$ and $e^{-1}$
gives mutually inverse maps of the actual face sets with each cardinality.
The resulting bijection preserves every $f_d$.
\end{proof}

The whole chain-incidence complex contains more data than the numeric count
vector. No theorem here states that the latter determines its isomorphism
class, reconstructs the original polytope, or is equal to its pre-subdivision
$f$-vector. For example a square boundary subdivides into eight vertices and
eight edges, while the full square subdivides into face counts $(9,16,8)$.
These examples clarify the two conventions; they are not the universal proof.

\section{Actual Face Families and Formal Scope}
Lean proves the above results for every poset and every finite family
$\mathcal A$ of \emph{actual subsets} of any ambient type $E$. It filters out
empty sets, uses actual subset inclusion, and constructs the order complex.
The boundary constructor additionally removes the specified whole face.
The flagness and $f$-vector chain formulas follow for both conventions. An
arbitrary input family need not itself be polytopal: the universal theorem
is stronger, and specializes to the actual $\mathcal F(Q)$ of a polytope.
No downward closure under arbitrary original \emph{vertex} subsets is required.
Thus a quadrilateral original face is one subdivision vertex, not a tetrahedron.

The standard combinatorial face-poset definition and the finiteness of the
actual face family of a bounded polytope identify the specialization in ordinary
mathematics. The Lean statements are explicitly universal finite-family
statements, not theorems about a separate geometric Polytope type. No barycenter
coordinates, affine independence, geometric gluing or coverage are advertised
as formally constructed. Such stronger geometric claims would require further
bridges and are not obtained merely by giving an arbitrary family a geometric name.

The core consists of 21 public theorems, plus the proved face and
transport equivalence definitions. Actual source replay and every public
theorem's compiled type and axiom footprint pass on Lean 4.33.1 with Mathlib
commit:
\begin{quote}\small\ttfamily
0df444a360eaa60ab8c11dca51a86af692955474.
\end{quote}
Only \texttt{propext}, \texttt{Classical.choice} and \texttt{Quot.sound} occur.
The portable package reproducer checks all nine physical Git dependencies,
their pinned origins/HEADs, tracked raw bytes and index/tree agreement before
and after execution. It cleans only this project's build outputs, actually
builds Main, replays the mathematical source and Main with warnings as errors,
and checks all 21 compiled theorem types and axiom footprints. Its finite face
incidence examples are regression checks, not the general proof. Installed Lean
and dependency build caches remain trust inputs, not an empty-machine bootstrap,
clean build of all Mathlib dependencies or formal geometric triangulation
theorem. The separate AI-agent review record does not imply human expertise,
organizer acceptance, first-solver priority or an award.

\begin{thebibliography}{9}
\bibitem{BW} F. Brenti and V. Welker, \emph{$f$-vectors of barycentric subdivisions}.
  Author-hosted manuscript, September 2007, especially p.~11.
\bibitem{OS} J. A. Olarte and F. Santos, \emph{Hypersimplicial subdivisions}.
  Selecta Mathematica \textbf{28} (2022), article~4 (online 2021).
  DOI: \texttt{10.1007/s00029-021-00722-6}, Section~2.
\end{thebibliography}
\end{document}
